% TEMPLATE-MILC-v3.tex - plantilla maestra para todas las guias MILC v3
% Ver scripts/generadores/build-guia-milc.py para la lista completa de
% claves esperadas. El script valida que no queden placeholders sin
% reemplazar antes de escribir el archivo final.

\documentclass[11pt,letterpaper]{article}
\usepackage[margin=1.75cm]{geometry}
\usepackage[table]{xcolor}
\usepackage{fontspec}
\usepackage[spanish]{babel}
\usepackage{tikz}
\usepackage[most]{tcolorbox}
\usepackage{tabularx}
\usepackage{array}
\usepackage{fancyhdr}
\usepackage{titlesec}
\usepackage{ragged2e}
\usepackage{enumitem}
\usepackage{microtype}

\setmainfont{Helvetica Neue}
\setsansfont{Helvetica Neue}
\setlength{\parindent}{0pt}
\setlength{\parskip}{6pt}
\hyphenpenalty=200
\exhyphenpenalty=200
\pretolerance=400
\tolerance=2000
\setlength{\emergencystretch}{1em}
\renewcommand{\arraystretch}{1.25}
\setlist{leftmargin=5mm,itemsep=1mm,topsep=1mm}
\newcolumntype{Y}{>{\RaggedRight\arraybackslash}X}

\definecolor{milcMagenta}{HTML}{E6007E}
\definecolor{milcVino}{HTML}{5A0038}
\definecolor{milcTurquesa}{HTML}{0093A5}
\definecolor{milcVerde}{HTML}{58A923}
\definecolor{milcMostaza}{HTML}{E5B400}
\definecolor{milcOcre}{HTML}{B86600}
\definecolor{milcNegro}{HTML}{241816}
\definecolor{milcGris}{HTML}{F7F7F4}
\definecolor{milcLinea}{HTML}{E8E2DD}
\definecolor{milcGradePrimary}{HTML}{84CC16}
\definecolor{milcGradeDark}{HTML}{4D7C0F}
\definecolor{milcGradeSoft}{HTML}{ECFCCB}
\definecolor{milcGradeAccent}{HTML}{CFE7FF}
\definecolor{milcGradeText}{HTML}{FFFFFF}
\color{milcNegro}

\pagestyle{fancy}
\fancyhf{}
\lhead{\small Tecnología e Informática · Grado 7}
\rhead{\small Guía 2 · MILC v3}
\cfoot{\small\thepage}
\renewcommand{\headrulewidth}{0pt}

\newtcolorbox{titlebox}[1]{
  enhanced,colback=#1,colframe=#1,arc=7mm,boxrule=0pt,
  left=7mm,right=7mm,top=3mm,bottom=3mm,
  before skip=8pt,after skip=8pt,
  fontupper=\sffamily\bfseries\Large\color{white}
}

\newtcolorbox{softbox}[3]{
  enhanced,colback=#2,colframe=#1,borderline west={4pt}{0pt}{#1},
  arc=2mm,boxrule=.4pt,left=4mm,right=4mm,top=3mm,bottom=3mm,
  before skip=6pt,after skip=8pt,title={#3},coltitle=white,
  colbacktitle=#1,fonttitle=\sffamily\bfseries,fontupper=\RaggedRight,
  attach boxed title to top left={xshift=3mm,yshift=-2mm},
  boxed title style={arc=2mm, boxrule=0pt}
}

\newtcolorbox{drawbox}[2]{
  enhanced,colback=white,colframe=#1,arc=4mm,boxrule=1pt,
  left=5mm,right=5mm,top=4mm,bottom=4mm,before skip=8pt,after skip=8pt,
  title={#2},coltitle=white,colbacktitle=#1,fonttitle=\sffamily\bfseries,
  fontupper=\RaggedRight,
  attach boxed title to top left={xshift=4mm,yshift=-2mm},
  boxed title style={arc=3mm, boxrule=0pt}
}

\newtcolorbox{infoband}[1]{
  enhanced,colback=#1!8,colframe=#1!50,arc=2mm,boxrule=.5pt,
  left=4mm,right=4mm,top=2mm,bottom=2mm,before skip=4pt,after skip=4pt,
  fontupper=\small\RaggedRight
}

% Puente narrativo entre secciones (contrato editorial Conéctate).
% Da continuidad: "Acabas de X. Ahora vas a Y porque Z."
% Visualmente: banda izquierda en color del grado, texto en itálica pequeña.
\newtcolorbox{puentebox}{
  enhanced,colback=milcGradeSoft,colframe=milcGradeDark,
  borderline west={3pt}{0pt}{milcGradeDark},
  arc=0mm,boxrule=0pt,left=5mm,right=4mm,top=2.5mm,bottom=2.5mm,
  before skip=4pt,after skip=8pt,
  fontupper=\itshape\small\color{milcNegro}\RaggedRight
}
\newcommand{\puente}[1]{%
  \begin{puentebox}%
    \textcolor{milcGradeDark}{\textbf{\textrightarrow}}\hspace{2mm}#1%
  \end{puentebox}%
}

\newcommand{\linead}[1]{\noindent\color{milcLinea}\rule{#1}{0.5pt}\color{milcNegro}}
\newcommand{\respuesta}[1]{\par\vspace{2mm}\foreach \n in {1,...,#1}{\noindent\color{milcLinea}\rule{\linewidth}{0.5pt}\par\vspace{5mm}}\color{milcNegro}}
\newcommand{\checkbox}{\textcolor{milcGradeDark}{\fbox{\phantom{x}}}\hspace{5pt}}

\newcommand{\phasecard}[4]{
\begin{tcolorbox}[enhanced,colback=#3,colframe=#2,arc=3mm,boxrule=.5pt,height=3.05cm,valign=center,halign=center,left=1.5mm,right=1.5mm,top=2mm,bottom=2mm]
{\sffamily\bfseries\fontsize{22}{24}\selectfont\textcolor{#2}{#1}}\par
{\sffamily\bfseries\small #4}
\end{tcolorbox}
}

\begin{document}
\thispagestyle{empty}
\begin{tikzpicture}[remember picture,overlay]
  \fill[white] (current page.south west) rectangle (current page.north east);
  \path[fill=milcGradePrimary,rounded corners=16mm]
    ([xshift=.55cm,yshift=-.55cm]current page.north west) rectangle
    ([xshift=-.55cm,yshift=3.65cm]current page.south east);

  \node[anchor=north west,text=milcGradeText] at ([xshift=2.0cm,yshift=-1.85cm]current page.north west)
    {\sffamily\bfseries\fontsize{14}{16}\selectfont GRADO};
  \node[anchor=north west,text=milcGradeText,opacity=.82] at ([xshift=2.0cm,yshift=-2.75cm]current page.north west)
    {\sffamily\bfseries\fontsize{9}{11}\selectfont SERIE GUÍAS · TECNOLOGÍA E INFORMÁTICA};

  \begin{scope}[shift={([xshift=-3.05cm,yshift=-2.55cm]current page.north east)}]
    \fill[white,opacity=.80] (-.62,-.52) rectangle (.62,.52);
    \draw[milcGradeText,line width=1pt] (-.62,-.52) rectangle (.62,.52);
    \fill[milcGradeDark] (-.42,-.38) rectangle (-.22,.06);
    \fill[milcGradeAccent] (-.10,-.38) rectangle (.10,.24);
    \fill[milcGradePrimary!60!black] (.22,-.38) rectangle (.42,.38);
  \end{scope}

  \node[anchor=west,text=milcGradeText] at ([xshift=1.9cm,yshift=-12.85cm]current page.north west)
    {\sffamily\bfseries\fontsize{124}{124}\selectfont 7};
  \draw[milcGradeText,line width=1.7pt]
    ([xshift=4.52cm,yshift=-11.68cm]current page.north west) circle (.13cm);

  \node[anchor=west,text=milcGradeText,text width=8.7cm,align=left] at ([xshift=10.35cm,yshift=-11.6cm]current page.north west)
    {
     {\sffamily\bfseries\fontsize{19}{22}\selectfont Microsoft\\365 ---\\las herramientas\\del taller\\digital}\\[4mm]
     {\sffamily\bfseries\fontsize{23}{26}\selectfont Guía 2-7-TIC}\\[2mm]
     {\sffamily\fontsize{13}{16}\selectfont Período 1 · Trabajo colaborativo en la nube}\\[5mm]
     {\sffamily\bfseries\fontsize{11}{13}\selectfont Institución Educativa Sor María Juliana}\\[1mm]
     {\sffamily\fontsize{10}{12}\selectfont Autor: Dr. Álvaro Cárdenas Orozco}
    };

  \path[fill=milcGradeSoft,rounded corners=7mm]
    ([xshift=1.35cm,yshift=.85cm]current page.south west) rectangle
    ([xshift=-1.35cm,yshift=4.25cm]current page.south east);
  \draw[milcGradeDark,line width=.65pt,rounded corners=7mm]
    ([xshift=1.35cm,yshift=.85cm]current page.south west) rectangle
    ([xshift=-1.35cm,yshift=4.25cm]current page.south east);
  \node[anchor=center,text=milcNegro,text width=17.0cm,align=center] at ([yshift=2.55cm]current page.south)
    {\sffamily\fontsize{10}{12}\selectfont Metodología MILC · Apertura ancestral · Pensamiento computacional · Triángulo Dussel-Estoicismo-Floridi\\
    \textcolor{milcGradeDark}{\bfseries Producto:} Un \textbf{mapa de las 8 aplicaciones principales de Microsoft 365} dibujado en una hoja del cuaderno, organizadas en \textbf{3 grupos}: \emph{(a) Documentos} (Word, Excel, PowerPoint, OneNote); \emph{(b) Comunicación} (Outlook, Teams); \emph{(c) Almacenamiento + más} (OneDrive, Forms). Cada app: \textbf{ícono dibujado, nombre, qué hace en 1 línea, ejemplo concreto de uso}. Más \textbf{una tabla comparativa} Microsoft 365 vs Google Workspace (Drive, Docs, Gmail, Meet) en 5 puntos.};
\end{tikzpicture}
\mbox{}
\newpage

\section*{Apertura: Microsoft 365 --- conocer las herramientas del taller digital}

\begin{softbox}{milcGradeDark}{milcGradeSoft}{Saber ancestral}
\textbf{Cuando llegabas a una minga grande en el campo de Cartago, lo primero era reconocer las herramientas.} En la finca cafetera del Valle del Cauca, alguien organizaba las herramientas de la minga al inicio: \textbf{canastas} para recoger el café, \textbf{lonas} para tender el grano al sol, \textbf{calabozos} (machetes pequeños) para podar, \textbf{ollas grandes} para la comida, \textbf{cucharas de madera} para servir, \textbf{cántaros} para llevar agua. Cada herramienta tenía su uso: \emph{canasta no servía para podar, machete no servía para servir sopa}. Doña Mercedes la maestra rural enseñaba a los niños que iban por primera vez: \emph{``Mire bien cada herramienta. Sepa cuál es cuál. No agarre cualquiera. La minga funciona si cada quien sabe qué herramienta usar''}. Hoy en tu taller digital pasa lo mismo: Microsoft 365 te ofrece \textbf{8 herramientas principales}. Cada una tiene su uso. Conocerlas todas te hace eficiente; usar la equivocada te complica el trabajo.
\end{softbox}

\begin{softbox}{milcTurquesa}{milcGris}{Saber contemporáneo}
\textbf{Microsoft 365} (antes llamado \emph{Office 365}) es un \textbf{paquete de aplicaciones} que Microsoft ofrece en la nube. \textbf{Tu colegio tiene licencia institucional}: te da una cuenta gratis con \texttt{tunombre@sormariajuliana.edu.co} (o similar) que abre las 8 aplicaciones desde cualquier dispositivo con internet. Las \textbf{8 aplicaciones principales} se agrupan en 3 familias. \textbf{Familia 1 (Documentos):} \emph{Word} (escribir), \emph{Excel} (cálculos y tablas), \emph{PowerPoint} (presentaciones), \emph{OneNote} (cuaderno digital). \textbf{Familia 2 (Comunicación):} \emph{Outlook} (correo), \emph{Teams} (chat, llamada, videollamada, reuniones). \textbf{Familia 3 (Almacenamiento + más):} \emph{OneDrive} (la nube donde se guardan tus archivos), \emph{Forms} (para crear encuestas y quizzes). Conocer las 8 y \textbf{cuándo usar cada una} te ahorra horas de mes y te hace eficiente como pocos compañeros lo serán.
\end{softbox}

\begin{softbox}{milcMostaza}{milcGris}{Pregunta-puente}
\textit{Si tu profe te pide \emph{``hacer una encuesta a 30 compañeros sobre qué materia les gusta más, recoger las respuestas en una tabla con gráficos, y presentarla a la clase''}, ¿\textbf{qué aplicación usarías para cada paso}? ¿Sabrías escogerlas o usarías Word para todo?}
\end{softbox}

\begin{softbox}{milcVerde}{milcGris}{Saber y saber hacer}
Al terminar la clase: (1) podrás \textbf{identificar} las 8 aplicaciones principales de Microsoft 365; (2) sabrás \textbf{explicar} qué hace cada una con 1 ejemplo concreto; (3) podrás \textbf{aplicar} la app correcta a 5 tareas reales; (4) habrás \textbf{comparado} Microsoft 365 con Google Workspace en 5 puntos.
\end{softbox}

\small
\begin{tabularx}{\textwidth}{>{\bfseries}p{3.0cm}Y}
\rowcolor{milcGradePrimary!12}Autor & Dr. Álvaro Cárdenas Orozco\\
DBA & Utiliza herramientas digitales para trabajar de manera colaborativa con otros (MEN, T\&I 7°)\\
Referentes & Saber ancestral de la minga · Microsoft 365 y OneDrive · Coautoría asincrónica\\
Producto final & Un \textbf{mapa de las 8 aplicaciones principales de Microsoft 365} dibujado en una hoja del cuaderno, organizadas en \textbf{3 grupos}: \emph{(a) Documentos} (Word, Excel, PowerPoint, OneNote); \emph{(b) Comunicación} (Outlook, Teams); \emph{(c) Almacenamiento + más} (OneDrive, Forms). Cada app: \textbf{ícono dibujado, nombre, qué hace en 1 línea, ejemplo concreto de uso}. Más \textbf{una tabla comparativa} Microsoft 365 vs Google Workspace (Drive, Docs, Gmail, Meet) en 5 puntos.\\
\end{tabularx}
\normalsize

\puente{Hoy haces 4 cosas: identificas las 8 apps, las agrupas en 3 familias, comparas con Google Workspace, y armas tu mapa visual.}

\begin{titlebox}{milcGradeDark}
Ruta MILC: así vamos a aprender
\end{titlebox}
\begin{tabularx}{\textwidth}{>{\centering\arraybackslash}X>{\centering\arraybackslash}X>{\centering\arraybackslash}X>{\centering\arraybackslash}X}
\phasecard{1}{milcVerde}{milcGris}{Escuta\\[-1mm]\small Observo, escucho y pregunto.} &
\phasecard{2}{milcTurquesa}{milcGris}{Sistematización\\[-1mm]\small Pensamiento computacional.} &
\phasecard{3}{milcMagenta}{milcGris}{Praxis\\[-1mm]\small Construyo y aplico.} &
\phasecard{4}{milcVino}{milcGris}{Evaluación\\[-1mm]\small Reflexión liberadora.}
\end{tabularx}


\puente{Empezamos con un test: ¿qué app usarías para 6 tareas distintas?}

\begin{titlebox}{milcVerde}
Fase 1 · Escuta
\end{titlebox}

\begin{softbox}{milcVerde}{milcGris}{Escena para escuchar}
\textbf{Actividad 1 · IDENTIFICA --- ¿Qué app usarías?} (10 min · individual). En el cuaderno, lee estas 6 tareas y escribe qué aplicación usarías para cada una (puedes adivinar; no se vale ver la sistematización). \textbf{Tarea 1:} Escribir un ensayo de 3 páginas sobre la independencia. \textbf{Tarea 2:} Crear una hoja con calificaciones de 30 estudiantes y calcular el promedio. \textbf{Tarea 3:} Hacer 20 diapositivas para una presentación oral. \textbf{Tarea 4:} Mandar correo formal al coordinador. \textbf{Tarea 5:} Reunión virtual con 5 compañeros para hacer trabajo en grupo. \textbf{Tarea 6:} Encuesta a toda la clase sobre cuál es la materia favorita.
\end{softbox}

\checkbox Leí las 6 tareas con calma.\\
\checkbox Escribí mi intuición de qué app usar en cada una.\\
\checkbox No me ayudé del compañero ni vi la sistematización.

\begin{infoband}{milcVerde}
\textbf{Lo que va al cuaderno (Actividad 1 · IDENTIFICA).} Título: «Actividad 1 --- 6 tareas y su app». \textbf{Formato:} lista numerada del 1 al 6 con tarea + app intuida. \textbf{Extensión:} 6 líneas. \textbf{Sabes que terminaste cuando} tienes una propuesta para cada una de las 6 tareas.
\end{infoband}

\puente{Después aprendes el nombre y función de las 8 apps, agrupadas en 3 familias.}

\begin{titlebox}{milcTurquesa}
Fase 2 · Sistematización con pensamiento computacional
\end{titlebox}

\begin{softbox}{milcTurquesa}{milcGris}{Cuatro pilares aplicados al tema}
Tu intuición te dio una idea. Ahora aprende las \textbf{8 aplicaciones principales} con sus familias, sus funciones y sus ejemplos exactos. Después podrás corregir tu primera Actividad si te equivocaste.
\end{softbox}

\small
\begin{tabularx}{\textwidth}{>{\bfseries\RaggedRight\arraybackslash}p{3.6cm}Y}
\rowcolor{milcTurquesa!90}\color{white}Pilar & \color{white}\bfseries Aplicación al tema\\
1. Descomposición & \textbf{Familia 1 --- Documentos (4 apps).} \textbf{Word}: escribir textos largos (ensayos, informes, cartas). Es el que ya conoces. \emph{Ícono: una W azul.} \textbf{Excel}: hacer cálculos y tablas con números. Ideal para presupuestos, calificaciones, listas con datos. \emph{Ícono: una X verde.} \textbf{PowerPoint}: hacer presentaciones con diapositivas. Útil para exponer un tema con apoyo visual. \emph{Ícono: una P naranja.} \textbf{OneNote}: cuaderno digital. Te permite tomar apuntes con texto, dibujos, imágenes, audios, todo en un solo lugar. \emph{Ícono: una N morada.}\\
2. Reconocimiento de patrones & \textbf{Familia 2 --- Comunicación (2 apps).} \textbf{Outlook}: correo electrónico institucional. Lo usas para enviar mensajes formales a profesores, coordinadores, padres. Tiene calendario integrado. \emph{Ícono: un sobre azul.} \textbf{Teams}: comunicación en tiempo real. Combina chat (como WhatsApp pero profesional), llamadas, videollamadas, reuniones de grupo, compartir pantalla. Si tu equipo va a hacer reunión virtual, usa Teams. \emph{Ícono: una T azul-morada.}\\
3. Abstracción & \textbf{Familia 3 --- Almacenamiento + más (2 apps).} \textbf{OneDrive}: la nube donde se guardan TODOS tus archivos. Cuando creas un Word o Excel en línea, va directo a OneDrive. Lo aprenderás en S3. \emph{Ícono: una nube azul.} \textbf{Forms}: para crear encuestas y quizzes. Recoge las respuestas automáticamente. Útil para investigaciones de aula. \emph{Ícono: un cuadro con líneas verdes/grises.}\\
4. Algoritmo conceptual & \textbf{Comparación Microsoft 365 vs Google Workspace.} Google también tiene su paquete: \emph{Docs} (= Word), \emph{Sheets} (= Excel), \emph{Slides} (= PowerPoint), \emph{Gmail} (= Outlook), \emph{Meet} (= Teams), \emph{Drive} (= OneDrive), \emph{Forms} (igual). \textbf{Diferencias:} \emph{Microsoft 365} suele tener \emph{más funciones profesionales}; \emph{Google Workspace} es \emph{más simple e intuitivo}, ideal para empezar. Ambos funcionan en la nube y permiten coautoría. \textbf{Pista:} aprender uno te facilita el otro, los conceptos son los mismos.\\
\end{tabularx}
\normalsize

\begin{drawbox}{milcTurquesa}{8 apps en 3 familias + atajos Google}
\textbf{Familia 1 (Documentos):} Word · Excel · PowerPoint · OneNote. \\[1mm]
\textbf{Familia 2 (Comunicación):} Outlook · Teams. \\[1mm]
\textbf{Familia 3 (Almacenamiento+):} OneDrive · Forms. \\[1mm]
\textbf{Equivalentes Google:} Docs · Sheets · Slides · (Keep) · Gmail · Meet · Drive · Forms.
\end{drawbox}

\begin{softbox}{milcMostaza}{milcGris}{Errores comunes y su corrección}
\textbf{Errores frecuentes al elegir app.} (1) \emph{Usar Word para todo} --- incluso para tablas con cálculos. Mal: usa Excel. (2) \emph{Hacer una presentación en Word con saltos de página} --- Mal: usa PowerPoint, está hecho para eso. (3) \emph{Mandar correo formal por WhatsApp} --- los correos profesionales van por Outlook, no por chat informal. (4) \emph{Hacer reunión virtual por WhatsApp con audio amateur} --- Teams o Google Meet hacen videollamada con compartir pantalla. (5) \emph{Hacer encuesta a mano y pasar resultados a Excel después} --- Forms recoge automáticamente, ahorrando horas.
\end{softbox}

\begin{infoband}{milcTurquesa}
\textbf{Lo que va al cuaderno (Actividad 2 · EXPLICA).} Título: «Actividad 2 --- 8 apps + 3 familias». \textbf{Formato:} tabla de 8 filas y 3 columnas. Columnas: \emph{App}, \emph{Familia}, \emph{Qué hace + ejemplo}. \textbf{Extensión:} 8 filas. \textbf{Sabes que terminaste cuando} podrías recomendarle la app correcta a un primo que te pregunte qué usar para una tarea.
\end{infoband}

\puente{Con todo claro, dibujas tu mapa de las 8 apps y la tabla comparativa.}

\begin{titlebox}{milcMagenta}
Fase 3 · Praxis con pensamiento computacional
\end{titlebox}

\begin{softbox}{milcMagenta}{milcGris}{Construyendo el producto con los cuatro pilares}
\textbf{Actividad 3 · CREA --- Mapa de Microsoft 365 + comparación con Google} (28 min · individual). Vas a hacer en una hoja del cuaderno un \textbf{mapa visual} de las 8 apps en 3 familias, más una \textbf{tabla comparativa} de Microsoft 365 vs Google Workspace.
\end{softbox}

\small
\begin{tabularx}{\textwidth}{>{\bfseries\RaggedRight\arraybackslash}p{3.6cm}Y}
\rowcolor{milcMagenta!90}\color{white}Pilar & \color{white}\bfseries Aplicación al producto\\
Descomposición & \textbf{Paso 1 --- Título y formato.} En la parte superior escribe: \textbf{``Mi mapa de Microsoft 365''}. Subraya. Pon la hoja en formato apaisado (horizontal) si puedes, para tener más espacio lateral.\\
Patrones & \textbf{Paso 2 --- Las 3 familias dibujadas.} Divide la hoja en \textbf{3 columnas verticales}. Encabezados: \emph{(1) DOCUMENTOS}, \emph{(2) COMUNICACIÓN}, \emph{(3) ALMACENAMIENTO + MÁS}. En la columna 1 dibuja 4 íconos (Word, Excel, PowerPoint, OneNote) en columna; en la 2 dibuja 2 íconos (Outlook, Teams); en la 3 dibuja 2 íconos (OneDrive, Forms).\\
Abstracción & \textbf{Paso 3 --- Para cada app, 3 datos.} Debajo o al lado de cada ícono escribe: \emph{nombre} (Word, etc.), \emph{qué hace en 1 línea}, \emph{1 ejemplo de uso} (\emph{``Escribir tarea de historia''}, \emph{``Calificaciones de 30''}, etc.). Sé conciso: no narres, solo datos.\\
Algoritmo de armado & \textbf{Paso 4 --- Tabla comparativa Microsoft 365 vs Google Workspace.} En la parte inferior de la hoja (o en otra hoja si te queda poco espacio), dibuja una tabla de \textbf{5 filas y 3 columnas}. Columnas: \emph{Función}, \emph{Microsoft 365}, \emph{Google Workspace}. Filas: \emph{(1) Escribir documentos --- Word vs Docs}; \emph{(2) Cálculos en tabla --- Excel vs Sheets}; \emph{(3) Presentaciones --- PowerPoint vs Slides}; \emph{(4) Correo --- Outlook vs Gmail}; \emph{(5) Reuniones virtuales --- Teams vs Meet}.\\
\end{tabularx}
\normalsize

\begin{drawbox}{milcMagenta}{Antes de mostrar tu mapa}
\checkbox El mapa tiene las 3 familias en columnas. \\[1mm]
\checkbox Las 8 apps tienen ícono dibujado, nombre, función y ejemplo. \\[1mm]
\checkbox La tabla comparativa Microsoft 365 vs Google tiene 5 filas. \\[1mm]
\checkbox Las 6 tareas de la Actividad 1 están corregidas. \\[1mm]
\checkbox La reflexión personal está en 2 líneas. \\[1mm]
\checkbox Está firmado con nombre y fecha.
\end{drawbox}

\begin{softbox}{milcMagenta}{milcGris}{Plantilla de trabajo}
\textbf{Estructura.} \textit{Arriba:} título + nombre. \textit{Centro:} 3 columnas (familias) con 4+2+2 íconos. Cada ícono: nombre, función, ejemplo. \textit{Abajo:} tabla comparativa 5x3. \textit{Final:} corrección Actividad 1 + reflexión 2 líneas + firma.
\end{softbox}

\begin{infoband}{milcMagenta}
\textbf{Lo que va al cuaderno (Actividad 3 · CREA).} Título: «Actividad 3 --- Mi mapa de Microsoft 365». \textbf{Formato:} mapa con 3 columnas + tabla comparativa + reflexión. \textbf{Extensión:} 1 hoja completa. \textbf{Sabes que terminaste cuando} podrías abrir la hoja en cualquier momento y reconocer qué hace cada app sin abrir el computador.
\end{infoband}

\puente{El producto es ese mapa + tabla comparativa firmados.}

\begin{titlebox}{milcMostaza}
Producto de la sesión
\end{titlebox}
\textbf{Mapa de Microsoft 365 con 3 familias + tabla vs Google · firmado}.
\begin{softbox}{milcMostaza}{milcGris}{Criterios de calidad}
(1) El mapa tiene las 8 apps en 3 familias correctamente agrupadas. (2) Cada app tiene ícono, nombre, función y ejemplo. (3) La tabla comparativa Microsoft 365 vs Google está completa. (4) Las correcciones de la Actividad 1 están hechas. (5) El mapa está firmado.
\end{softbox}

\puente{Antes de salir, verifica que cada app tiene su ícono, nombre, función y ejemplo.}

\begin{titlebox}{milcVino}
Fase 4 · Evaluación liberadora — cinco dimensiones
\end{titlebox}

\begin{softbox}{milcVino}{milcGris}{Sentido de la evaluación}
Evaluar no es solo revisar una nota. Es reconocer qué aprendí, qué cambió en mí, qué emociones manejé, cómo me relaciono con la ciudad y la comunidad, y qué le dejaré a quien viene detrás.
\end{softbox}

\textbf{Mi reflexión liberadora en cinco dimensiones.} Completa cada frase iniciada:

\small
\begin{tabularx}{\textwidth}{>{\bfseries\RaggedRight\arraybackslash}p{3.4cm}Y>{\RaggedRight\arraybackslash}p{4.6cm}}
\rowcolor{milcVino}\color{white}Dimensión & \color{white}\bfseries Mi respuesta (frame starter) & \color{white}\bfseries Algo que descubrí\\
1. Personal & Lo que más cambió en mí fue \linead{6.5cm} & \linead{4.2cm}\\
2. Emocional & La emoción más fuerte fue \linead{2.5cm} y la regulé \linead{3cm} & \linead{4.2cm}\\
3. Ciudadana & En democracia esto importa porque \linead{6cm} & \linead{4.2cm}\\
4. Local (Cartago) & En mi barrio sería útil para \linead{6.5cm} & \linead{4.2cm}\\
5. Intergeneracional & A mi abuela/abuelo le diría \linead{6.5cm} & \linead{4.2cm}\\
\end{tabularx}
\normalsize

\begin{infoband}{milcVino}
\textbf{Para la próxima sustentación.} Vuelve a esta tabla en seis meses. Verás que las respuestas cambian — no porque lo aprendido cambie, sino porque tú cambias. La evaluación liberadora es una conversación contigo a través del tiempo.
\end{infoband}


\puente{Cerramos con 3 voces que te ayudan a entender por qué conocer las herramientas es la primera ventaja.}

\begin{titlebox}{milcGradeDark}
Triángulo de pensamiento · cierre filosófico
\end{titlebox}

\begin{softbox}{milcVino}{milcGris}{Dussel · lente del nosotros}
\textit{````Quien conoce las herramientas, manda; quien no las conoce, obedece a quien sí.''''}

Hoy hay personas que viven en el siglo XXI \textbf{sin conocer las herramientas digitales básicas}. Hacen su trabajo a mano, mandan documentos en papel, hacen reuniones presenciales para todo. \emph{No están mal, pero pierden tiempo}. Aprender las 8 apps de Microsoft 365 (o sus equivalentes Google) te coloca en el grupo de los que \emph{eligen} sus herramientas, no de los que dependen de que alguien más decida. Es alfabetización del siglo XXI.

\textbf{¿Estoy aprendiendo las herramientas que me darán libertad en mi vida adulta?}
\end{softbox}
\linead{\linewidth}\\[1mm]
\linead{\linewidth}

\begin{softbox}{milcMostaza}{milcGris}{Estoicismo · Marco Aurelio · cuidado interior}
\textit{````Cada herramienta tiene su uso. Forzarla a otro le quita eficacia y dignidad. Conoce primero su naturaleza.''''}

Usar Word para hacer una tabla con cálculos es como usar un martillo para apretar tornillos: funciona mal y daña la herramienta. Cada app de Microsoft 365 tiene un uso para el que está diseñada. \textbf{Aprender ese uso es respetar la herramienta}, y a la vez ahorrarse problemas. Es disciplina de profesional, no perfeccionismo.

\textbf{¿Estoy usando cada herramienta para lo que sirve, o forzando una a hacer trabajo de otra?}
\end{softbox}
\linead{\linewidth}\\[1mm]
\linead{\linewidth}

\begin{softbox}{milcTurquesa}{milcGris}{Floridi · lente de la infoesfera}
\textit{````El paquete de herramientas que dominas define qué decisiones puedes tomar. Más herramientas, más libertad.''''}

Si solo conoces Word, vas a hacer todo en Word. Si conoces las 8 apps, \textbf{eliges la mejor para cada situación}. Esa elección te da \emph{ventaja sobre quien no las conoce}: más rápido, mejor resultado, menos esfuerzo. \emph{Más herramientas en tu caja = más libertad para escoger}. Por eso vale la pena la inversión inicial de aprender las 8.

\textbf{¿Qué app del paquete me da más curiosidad? ¿Por qué no me animo a explorarla?}
\end{softbox}
\linead{\linewidth}\\[1mm]
\linead{\linewidth}

\begin{titlebox}{milcVino}
Compromiso final
\end{titlebox}

\small
\begin{tabularx}{\textwidth}{>{\RaggedRight\arraybackslash}p{3.0cm}YYY}
\rowcolor{milcVino}\color{white}\bfseries Criterio & \color{white}\bfseries Lo logré & \color{white}\bfseries En proceso & \color{white}\bfseries Necesito apoyo\\
Comprensión del tema & Explico ideas centrales con mis palabras. & Reconozco ideas pero falta precisión. & Necesito repasar conceptos base.\\
Pensamiento computacional & Apliqué los 4 pilares al diseñar mi producto. & Apliqué algunos pilares. & Necesito releer la fase 2.\\
Aplicación y producto & Mi producto cumple los criterios y se puede revisar. & Producto funcional con ajustes pendientes. & Aún en construcción.\\
Reflexión 5 dimensiones & Respondí cada dimensión con honestidad. & Respondí algunas con detalle. & Mis respuestas son muy generales.\\
Triángulo de pensamiento & Conecté las 3 voces con mi trabajo real. & Conecté al menos una. & No relaciono las voces con mi proceso.\\
\end{tabularx}
\normalsize

\textbf{Hoy entendí que:}\\[1mm]
\linead{\linewidth}\\[1mm]
\linead{\linewidth}

\textbf{Mi compromiso para la próxima sesión es:}\\[1mm]
\linead{\linewidth}\\[1mm]
\linead{\linewidth}

\begin{softbox}{milcMostaza}{milcGris}{Cita de despedida · Marco Aurelio}
\textit{``El obstáculo en el camino se vuelve el camino. Lo que estorba al camino se vuelve camino.''}\\[1mm]
Cada error de hoy, bien observado, se convierte en pieza del camino que pisarás mañana.
\end{softbox}

\small
\begin{tabularx}{\textwidth}{>{\bfseries}p{3.6cm}Y}
\rowcolor{milcGradeDark!12}Nombre completo: & \linead{10cm}\\
Fecha: & \linead{4cm} \quad Curso/grupo: \linead{4cm}\\
Mi firma: & \linead{9cm}\\
\end{tabularx}
\normalsize

\begin{infoband}{milcGradeDark}
\textbf{Para la próxima sesión.} Esta semana, abre Microsoft 365 en tu cuenta institucional (\texttt{office.com}) y \emph{solo mira} las 8 aplicaciones disponibles. No tienes que crear nada todavía. Solo familiarizarse. La próxima clase activamos \textbf{OneDrive}, tu nube personal, y comenzamos a guardar archivos en ella.
\end{infoband}

\end{document}
